% from ega4-18-fr.tex
\oldpage[IV-4]{145}

(On notera que l'hypothèse de $d)$ est vérifiée si $f$ est \emph{formellement non ramifié} au point $x$ (17.4.1.2).)

Nous laissons au lecteur la démonstration, qui est substantiellement la même que celle de (18.5.18).

\begin{lemma}[18.8.4]
\label{IV.18.8.4}
Soient $A$, $A'$ deux anneaux locaux, $\varphi:A\to A'$ un homomorphisme local faisant de $A'$ une $A$-algèbre essentiellement étale (18.6.1). Alors, pour tout anneau strictement local $B$, tout homomorphisme local $\psi:A\to B$ et tout homomorphisme de $k(A)$-algèbres $\alpha:k(A')\to k(B)$, il existe un et un seul homomorphisme local $\psi':A'\to B$ tel que $\psi=\psi'\circ\varphi$ et que $\alpha$ soit égal à l'homomorphisme $\overline{\psi'}$ déduit de $\psi'$ par passage aux quotients.
\end{lemma}

Avec les notations de (18.6.2), l'homomorphisme $\alpha$ correspond à un $k(s)$-morphisme $\Spec(k(y))\to\Spec(k(x))$, donc à un point $x'$ bien déterminé dans $X'$ au-dessus de $y$. L'existence de la $Y$-section $f'$ de $X'$ telle que $f'(y)=x'$ résulte donc de (18.8.1, $c)$) puisque $f'$ est étale et $B$ strictement local, et son unicité résulte de ce que $f'$ est séparé et de (17.4.9).

\begin{lemma}[18.8.5]
\label{IV.18.8.5}
Soient $A$ un anneau local, $A_1$, $A_2$ deux $A$-algèbres locales essentiellement étales, $K$ un corps, extension de $k(A)$.
\begin{enumerate}
  \item[(i)] Pour tout $k(A)$-homomorphisme $\gamma:k(A_1)\to k(A_2)$, il existe au plus un $A$-homomorphisme (nécessairement local) $\psi:A_1\to A_2$ tel que $\gamma$ soit l'homomorphisme $\overline\psi$ déduit de $\psi$ par passage aux quotients.
  \item[(ii)] Pour tout couple de $A$-homomorphismes locaux $\beta_1:A_1\to K$, $\beta_2:A_2\to K$ il existe une $A$-algèbre locale essentiellement étale $A_3$, deux $A$-homomorphismes $\varphi_1:A_1\to A_3$, $\varphi_2:A_2\to A_3$ et un $A$-homomorphisme $\beta:A_3\to K$ tels que $\beta_1=\beta\circ\varphi_1$ et $\beta_2=\beta\circ\varphi_2$.
\end{enumerate}
\end{lemma}

Avec les notations de (18.6.3), les homomorphismes $\beta_1$, $\beta_2$ dans (ii) correspondent à deux $S$-morphismes $\Spec(K)\to X_1$, $\Spec(K)\to X_2$, d'images respectives $x_1$, $x_2$; on en déduit un $S$-morphisme bien déterminé $\Spec(K)\to X_3$ (I, 3.2.1) d'image $x_3$ au-dessus de $x_1$ et $x_2$; la $A$-algèbre $A_3=\sh{O}_{X_3,x_3}$ et l'homomorphisme $A_3\to K$ correspondant à ces données répondent aux conditions de (ii). D'autre part, la donnée de $\gamma$ dans (i) correspond à un $S$-morphisme $\Spec(k(x_2))\to\Spec(k(x_1))$ ou encore à un point $x'_3$ de $X'_3$ au-dessus de $x_2$; l'unicité de $\psi$ résulte de l'unicité d'une $X_2$-section de $X'_3$ passant par $x'_3$ (17.4.9).

\begin{env}[18.8.6]
\label{IV.18.8.6}
Soient $A$ un anneau local, $\mathfrak m$ son idéal maximal, $k$ son corps résiduel, $\Omega$ une clôture séparable de $k$. Considérons l'ensemble $\mathfrak E$ de $A$-algèbres essentiellement étales défini dans (18.6.1), et désignons par $L$ l'ensemble des $A$-homomorphismes locaux $A'\to\Omega$, où $A'$ parcourt $\mathfrak E$; on notera qu'un tel homomorphisme se factorise en
\[
  \lambda:A'\longrightarrow k(A')\xrightarrow{\omega_\lambda}\Omega,
\]
où $\omega_\lambda:k(A')\to\Omega$ est un $k$-homomorphisme (puisque $\mathfrak mA'$ est l'idéal maximal de $A'$); inversement la donnée d'un tel homomorphisme $\omega_\lambda$ détermine uniquement un homomorphisme $\lambda\in L$. Pour tout $\lambda\in L$, nous désignerons par $A_\lambda$ la $A$-algèbre essentiellement étale appartenant à $\mathfrak E$ sur laquelle $\lambda$ est défini. L'ensemble $L$ est \emph{préordonné} par la relation « il existe un $A$-homomorphisme $\varphi_{\mu\lambda}:A_\lambda\to A_\mu$ tel que $\lambda=\mu\circ\varphi_{\mu\lambda}$ », et il résulte de (18.8.5, (i)) que cet homomorphisme $\varphi_{\mu\lambda}$ est \emph{unique}. En outre, $L$ est \emph{filtrant croissant} pour la relation de préordre précédente (que l'on note encore $\lambda\leq\mu$), en vertu

